%==============================================================
%  lecons/algebre/nombres_complexes.tex — Nombres complexes
%==============================================================

\lecontitle{Nombres complexes}{Structures algébriques — Arithmétique de \texorpdfstring{$\mathbb{C}$}{C}}

%=============================================================
%  § 1 — FORMES ET OPÉRATIONS ÉLÉMENTAIRES
%=============================================================
\secnum{navyblue}{1}{Représentations et opérations élémentaires}

\begin{tcolorbox}[thmbox]
  \textbf{Forme algébrique.}\enspace%
  \index{nombres complexes!forme algébrique}
  Tout $z \in \mathbb{C}$ s'écrit de façon unique
  \[
    z = a + ib, \quad a = \mathrm{Re}(z) \in \mathbb{R},\quad b = \mathrm{Im}(z) \in \mathbb{R}.
  \]
  Le \textbf{conjugué}\index{nombres complexes!conjugué} et le
  \textbf{module}\index{nombres complexes!module} sont :
  \[
    \bar{z} = a - ib, \qquad |z| = \sqrt{a^2 + b^2} = \sqrt{z\bar{z}}.
  \]
\end{tcolorbox}

\rubriq{Propriétés du conjugué et du module}

Les règles suivantes découlent directement des définitions. Pour tous
$z, w \in \mathbb{C}$ :
\[
  \overline{z + w} = \bar{z} + \bar{w}, \quad
  \overline{zw} = \bar{z}\,\bar{w}, \quad
  \overline{\bar{z}} = z, \quad
  z + \bar{z} = 2\,\mathrm{Re}(z), \quad
  z - \bar{z} = 2i\,\mathrm{Im}(z).
\]
Pour le module :
\[
  |zw| = |z|\,|w|, \quad
  \left|\frac{z}{w}\right| = \frac{|z|}{|w|} \;(w \neq 0), \quad
  |\bar{z}| = |z|, \quad
  |z + w| \leq |z| + |w|.
\]

\rubriq{Inverse et division}
Pour $z = a + ib \neq 0$ :
\[
  z^{-1} = \frac{\bar{z}}{|z|^2} = \frac{a - ib}{a^2 + b^2}.
\]
La division se ramène à une multiplication par le conjugué du dénominateur :
$\dfrac{z}{w} = \dfrac{z\bar{w}}{|w|^2}$.

\decsep{}

%=============================================================
%  § 2 — FORME TRIGONOMÉTRIQUE ET FORMULE D'EULER
%=============================================================
\secnum{navyblue}{2}{Forme trigonométrique et formule d'Euler}

\begin{tcolorbox}[thmbox]
  \textbf{Forme trigonométrique.}\enspace%
  \index{nombres complexes!forme trigonométrique}\index{argument}
  Tout $z \neq 0$ s'écrit
  \[
    z = |z|(\cos\theta + i\sin\theta), \quad \theta \in \mathbb{R}.
  \]
  Le réel $\theta$ est un \textbf{argument}\index{argument} de $z$, noté
  $\theta \equiv \arg(z) \pmod{2\pi}$. L'argument principal est choisi dans
  $]-\pi, \pi]$.

  \medskip
  \textbf{Formule d'Euler.}\enspace%
  \index{formule d'Euler}
  Pour tout $\theta \in \mathbb{R}$ :
  \[
    e^{i\theta} = \cos\theta + i\sin\theta.
  \]
  La forme exponentielle est donc $z = |z|\,e^{i\arg z}$.
\end{tcolorbox}

\rubriq{Multiplication en forme polaire}

\begin{mdframed}[style=proofbar]
  Si $z = r e^{i\theta}$ et $w = s e^{i\varphi}$, alors
  \[
    zw = rs\, e^{i(\theta + \varphi)}.
  \]
  \textit{Règles pour l'argument :}
  $\arg(zw) \equiv \arg z + \arg w \pmod{2\pi}$,
  $\arg(z^{-1}) \equiv -\arg z$,
  $\arg(\bar{z}) \equiv -\arg z$.

  \textit{Conséquence géométrique :} multiplier par $e^{i\alpha}$ est une
  \textbf{rotation}\index{rotation (multiplication complexe)} de centre $O$ et d'angle $\alpha$ dans le plan complexe.
  Multiplier par $r > 0$ est une homothétie de rapport $r$. \hfill$\blacksquare$
\end{mdframed}

\rubriq{Formules de De Moivre et de linéarisation}

\begin{tcolorbox}[thmbox]
  \textbf{Formule de De Moivre.}\enspace%
  \index{De Moivre (formule de)}
  Pour tout $n \in \mathbb{Z}$ et $\theta \in \mathbb{R}$ :
  \[
    (\cos\theta + i\sin\theta)^n = \cos(n\theta) + i\sin(n\theta).
  \]

  \medskip
  \textbf{Formules d'Euler inverses.}\enspace%
  \[
    \cos\theta = \frac{e^{i\theta} + e^{-i\theta}}{2}, \qquad
    \sin\theta = \frac{e^{i\theta} - e^{-i\theta}}{2i}.
  \]
  Elles permettent de \textbf{linéariser}\index{linéarisation trigonométrique} $\cos^n\theta$ et $\sin^n\theta$
  (exprimer une puissance en combinaison de $\cos(k\theta)$, $\sin(k\theta)$).
\end{tcolorbox}

\begin{mdframed}[style=proofbar]
  \textit{Preuve de De Moivre.} Par récurrence sur $n \geq 0$ en utilisant
  $e^{i(n+1)\theta} = e^{in\theta} \cdot e^{i\theta}$. Le cas $n < 0$ suit
  par passage à l'inverse. \hfill$\blacksquare$
\end{mdframed}

\decsep{}

%=============================================================
%  § 3 — RACINES N-IÈMES
%=============================================================
\secnum{navyblue}{3}{Racines $n$-ièmes}

\begin{tcolorbox}[thmbox]
  \textbf{Théorème (racines $n$-ièmes).}\enspace%
  \index{racines n-ièmes}
  Soit $n \geq 1$ et $w = \rho e^{i\varphi} \in \mathbb{C}^*$. L'équation
  $z^n = w$ admet exactement $n$ solutions dans $\mathbb{C}$ :
  \[
    z_k = \rho^{1/n}\, \exp\!\left(i\,\frac{\varphi + 2k\pi}{n}\right),
    \quad k = 0, 1, \ldots, n-1.
  \]
  Ces solutions sont les sommets d'un polygone régulier à $n$ côtés inscrit
  dans le cercle de rayon $\rho^{1/n}$, équidistants d'angle $2\pi/n$.

  \medskip
  \textbf{Racines de l'unité.}\enspace%
  \index{racines de l'unité}
  Pour $w = 1$ ($\rho = 1$, $\varphi = 0$) :
  \[
    \omega_k = e^{2ik\pi/n}, \quad k = 0, \ldots, n-1.
  \]
  Ces racines vérifient $\displaystyle\sum_{k=0}^{n-1} \omega_k^j = 0$
  pour $j \not\equiv 0 \pmod{n}$.
\end{tcolorbox}

\begin{mdframed}[style=proofbar]
  \textit{Existence.} On vérifie $z_k^n = \rho\,e^{i(\varphi+2k\pi)} = \rho e^{i\varphi} = w$.

  \textit{Unicité (exactement $n$ solutions).} Si $z^n = w$, on écrit $z = r e^{i\alpha}$.
  Alors $r^n = \rho$ donne $r = \rho^{1/n}$ (unique réel positif), et
  $n\alpha \equiv \varphi \pmod{2\pi}$, ce qui donne exactement $n$ valeurs
  $\alpha \in [0, 2\pi[$.

  \textit{Somme des racines.} Pour $j \not\equiv 0 \pmod{n}$, la somme est
  une série géométrique :
  $\sum_{k=0}^{n-1} \omega_k^j = \sum_{k=0}^{n-1} (e^{2i\pi j/n})^k
  = \dfrac{1 - e^{2i\pi j}}{1 - e^{2i\pi j/n}} = 0$.
  \hfill$\blacksquare$
\end{mdframed}

\decsep{}

%=============================================================
%  § 4 — GÉOMÉTRIE DU PLAN COMPLEXE
%=============================================================
\secnum{navyblue}{4}{Géométrie du plan complexe}

\begin{tcolorbox}[thmbox]
  \textbf{Transformations géométriques.}\enspace%
  \index{plan complexe!transformations}
  On identifie $\mathbb{C}$ et $\mathbb{R}^2$ via $z = x + iy \leftrightarrow (x,y)$.
  \begin{itemize}
    \item \textbf{Translation} de vecteur $b$ : $z \mapsto z + b$.
    \item \textbf{Homothétie} de centre $O$, rapport $r > 0$ : $z \mapsto rz$.
    \item \textbf{Rotation} de centre $O$, angle $\alpha$ : $z \mapsto e^{i\alpha} z$.
    \item \textbf{Similitude directe} (générale) : $z \mapsto az + b$,
          $a \in \mathbb{C}^*$, $b \in \mathbb{C}$.
    \item \textbf{Réflexion} d'axe réel : $z \mapsto \bar{z}$.
    \item \textbf{Inversion} de centre $O$, puissance $1$ : $z \mapsto 1/\bar{z}$.
  \end{itemize}
\end{tcolorbox}

\rubriq{Alignement et cocyclicité}

\begin{mdframed}[style=proofbar]
  Trois points $A$, $B$, $C$ d'affixes $a$, $b$, $c$ sont
  \textbf{alignés}\index{alignement (plan complexe)} si et seulement si
  \[
    \frac{c - a}{b - a} \in \mathbb{R}.
  \]
  Quatre points sont \textbf{cocycliques ou alignés}\index{cocyclicité} si et
  seulement si leur birapport\index{birapport}
  $[z_1, z_2; z_3, z_4] = \dfrac{(z_3-z_1)(z_4-z_2)}{(z_3-z_2)(z_4-z_1)}$
  est réel. \hfill$\blacksquare$
\end{mdframed}

\decsep{}

%=============================================================
%  § 5 — EXEMPLES, PIÈGES ET EXERCICES
%=============================================================
\secnum{exgreen}{5}{Exemples, pièges et exercices}

\noindent{\bfseries\color{navyblue} Exemples.}

\begin{mdframed}[style=exbar]
  \textbf{1. Linéarisation de $\cos^3\theta$.}\enspace%
  En développant $\left(\dfrac{e^{i\theta}+e^{-i\theta}}{2}\right)^3$ par le
  binôme de Newton :
  \[
    \cos^3\theta = \frac{1}{4}\cos(3\theta) + \frac{3}{4}\cos\theta.
  \]

  \smallskip\noindent
  \textbf{2. Racines cubiques de $-8$.}\enspace%
  $w = -8 = 8e^{i\pi}$, $n = 3$ : $z_k = 2\exp\!\bigl(i\dfrac{\pi + 2k\pi}{3}\bigr)$.
  Cela donne $z_0 = 2e^{i\pi/3} = 1 + i\sqrt{3}$, $z_1 = -2$,
  $z_2 = 1 - i\sqrt{3}$.

  \smallskip\noindent
  \textbf{3. Somme trigonométrique.}\enspace%
  $\displaystyle S_n = \sum_{k=0}^{n} \cos(k\theta) = \mathrm{Re}\!\sum_{k=0}^n e^{ik\theta}
  = \mathrm{Re}\,\frac{1 - e^{i(n+1)\theta}}{1 - e^{i\theta}}$.
  On factorise pour obtenir la formule fermée :
  \[
    S_n = \frac{\sin\!\bigl(\tfrac{(n+1)\theta}{2}\bigr)}
               {\sin\!\bigl(\tfrac{\theta}{2}\bigr)}\,
          \cos\!\left(\frac{n\theta}{2}\right), \quad \theta \notin 2\pi\mathbb{Z}.
  \]

  \smallskip\noindent
  \textbf{4. Identité d'Euler.}\enspace%
  $e^{i\pi} + 1 = 0$ : les cinq constantes fondamentales liées en une seule
  égalité.
\end{mdframed}

\vspace{6pt}
\noindent{\bfseries\color{counterred} Pièges.}

\begin{mdframed}[style=cexbar]
  \textbf{Argument non additif sur $\mathbb{R}$.}\enspace%
  $\arg(zw) = \arg z + \arg w$ n'est valide que modulo $2\pi$. L'argument
  principal peut présenter un saut de $\pm 2\pi$ : $\arg(-1 \cdot -1) =
  \arg(1) = 0$, mais $\arg(-1) + \arg(-1) = -\pi + (-\pi) = -2\pi \neq 0$.

  \smallskip\noindent
  \textbf{$\sqrt{z^2} \neq z$ en général.}\enspace%
  On ne peut pas définir une racine carrée globale continue sur $\mathbb{C}^*$
  (le logarithme complexe est multiforme). Par exemple,
  $\sqrt{(-i)^2} = \sqrt{-1} = i \neq -i$.

  \smallskip\noindent
  \textbf{Inégalité triangulaire : égalité.}\enspace%
  $|z + w| = |z| + |w|$ si et seulement si $\arg z \equiv \arg w \pmod{2\pi}$
  (même argument, c.-à-d.\ $w = \lambda z$ pour un $\lambda \geq 0$), pas
  simplement si $z$ et $w$ sont positifs.
\end{mdframed}

\vspace{6pt}
\noindent{\bfseries\color{goldaccent} Exercices.}

\begin{mdframed}[style=exobar]
  \begin{enumerate}
    \item Linéariser $\sin^4\theta$ et $\cos^2\theta\sin^2\theta$ à l'aide des
          formules d'Euler.

    \item Calculer les racines cinquièmes de l'unité. En déduire les valeurs
          exactes de $\cos(2\pi/5)$ et $\cos(4\pi/5)$.

    \item Soit $z_0 \in \mathbb{C}$, $|z_0| < 1$. Montrer que la transformation
          de Möbius\index{transformation de Möbius}
          $f(z) = \dfrac{z - z_0}{1 - \bar{z}_0 z}$
          est une bijection du disque unité ouvert sur lui-même, et que $|f(z)| = 1$
          si $|z| = 1$.

    \item Résoudre dans $\mathbb{C}$ l'équation $z^4 + 2z^2 + 4 = 0$.
          (\textit{Indication :} discriminant de l'équation en $z^2$.)

    \item Soit $P(z) = z^n + a_{n-1}z^{n-1} + \cdots + a_0 \in \mathbb{R}[X]$.
          Montrer que si $z_0$ est racine de $P$, alors $\bar{z}_0$ l'est aussi,
          et que les racines non réelles viennent par paires conjuguées.
          En déduire que tout polynôme réel de degré impair admet au moins une
          racine réelle.

    \item (Inégalité de Ptolémée) Montrer que pour tous $z_1, z_2, z_3 \in \mathbb{C}$ :
          \[
            |z_1 - z_3| \cdot |z_2| \;\leq\;
            |z_1 - z_2| \cdot |z_3| + |z_2 - z_3| \cdot |z_1|.
          \]
          (\textit{Indication :} écrire $z_1 - z_3 = (z_1 - z_2) + (z_2 - z_3)$
          et factoriser habilement.)
  \end{enumerate}
\end{mdframed}

\leconfooter{R.\ Descartes (1637) — L.\ Euler (1748) — A.\ de Moivre (1730) — C.\ F.\ Gauss (1799)}
